% ======================= 导 论 =======================
\chapter*{导\quad 论}
\phantomsection
\addcontentsline{toc}{chapter}{导\quad 论}
\markboth{导论}{导论}

\noindent
关于「月之都」这四个字，通行的理解是一座宫殿：琉璃为瓦，白玉为阶，住着一位不能回头的公主。这个理解有一半是错的，错的那一半不在宫殿，在「一座」。

它不止一座。月球正面静海西缘有一处设施，编号 MARE-TRANQ-07，中文材料里写作广寒宫；月球背面另有一座城，那才是「月之都」这三个字原本指的东西。两者同属一城，前者是这座城伸出正面的一角，后者是主体。第 I 卷记录的是那一角留下的三十七片残简，第 VII 卷记录的是那一角里的人的十二万后代如何改写了中国的历史。这一本写主体。

本书要证明的命题只有一条：月之都的一切制度，都可以从「把变化从自己身上删掉」这一件事推出来；而它的每一次失败，也都是这一件事的后果。

这条命题可以拆成三个可以分别检验的部分。其一，维持不变的成本高于变化本身的成本；其二，不变的人无法更新轻重缓急，因此无法纠错；其三，删掉变化的人，最终还是被变化处理掉。

\section*{四类材料}

本书用四类材料。

影像。一九五九年以后，人类对月球背面拍摄了数十次，分辨率从一公里提高到半米。这类材料最硬，也最容易被忽略，理由见第一章。

档案。`ФВ` 系列卷宗，一九六三年以后由苏联科学院远东委员会搜集，一九九二年移交 VFS 基金会。本书用到的编号集中在 1965 至 1972 年之间，凡引必给叶次。

文本。三十七片残简及其语译见第 I 卷，本书只引简次与结论，不重录释文。中国方面的传世文献（《左传》《淮南子》《史记》《竹取物语》）只用于比对，不用于立论。

地层。贝加尔湖 BDP-96 岩芯与月面热流数据，用于定时。

四类材料的分量并不相等。影像与地层可以复核，档案只能复核「有没有这一份」，文本则全凭释读。本书在每一章的末尾都会说明，该章的结论有多少重量压在可以复核的那一类上。

\section*{方法与限度}

笔者的办法是先摆器物，再谈它意味着什么。

这与训练有关。笔者没有受过史学训练，九〇年代受聘于 VFS 异常文献部做日本联络员，做的是收发与转写，不是研究。二〇二四年读到第 I 卷的校样之后才开始动手写这一本，写的时候依照的是三卷共用的那份体例：一条材料、一句推论、一处自我拆台。做惯了手艺的人，信不过没有工序的结论。

限度也要说清楚。笔者没有去过月面，见过的只有照片、一卷缩微胶片和巴兰尼科夫教授书里那几张重绘图。凡本书讲「背面长什么样」的地方，读者都可以先打个折。

\section*{全书的结构}

十章分四段。

第一、二章讲这座城是什么：它的形状，以及它为什么不动。

第三至第七章讲它的四次失败：造出来的手失去了控制；规程在两万年间没有改过；十部历被废掉九部；一台拆东西的机器被一个不识字的人用了；最后是月兔之乱。

第八、九章讲它的终结与余波：撤离、落地、渡海。

第十章讲名字，以及一个文明如何在别人的书里变成神话。

结语是笔者自己的一段交代。那一段用的是「我」，其余各章一律用「笔者」。

\begin{flushright}
{\small 太田顺也}\\[1mm]
{\footnotesize\color{pencil} 二〇二五年一月}
\end{flushright}

\clearpage
